\documentclass[11pt]{article}

\usepackage[utf8]{inputenc}
\usepackage[T1]{fontenc}
\usepackage{lmodern}
\usepackage{geometry}
\usepackage{amsmath,amssymb,amsfonts,amsthm,mathtools}
\usepackage{bm}
\usepackage{enumitem}
\usepackage{microtype}
\usepackage{hyperref}
\usepackage{titlesec}
\usepackage{booktabs}
\usepackage{array}
\usepackage{setspace}
\usepackage{mathrsfs}
\usepackage{physics}

\geometry{margin=1in}
\hypersetup{colorlinks=true, linkcolor=black, urlcolor=blue, citecolor=black}
\setlist[enumerate]{topsep=0.4em,itemsep=0.35em}
\setlist[itemize]{topsep=0.3em,itemsep=0.25em}
\titleformat{\section}{\large\bfseries}{\thesection.}{0.5em}{}
\titleformat{\subsection}{\normalsize\bfseries}{\thesubsection.}{0.5em}{}
\setstretch{1.06}

\newcommand{\Aether}{\AE{}ther}
\newcommand{\Aflow}{\AE{}ther-flow}
\newcommand{\TheoryName}{\Aflow{} Interpretation of Relativity}
\newcommand{\TheoryProgram}{\Aether{} / \Aflow{} framework}

\newcommand{\CompletedMetric}{g^{\sharp}}
\newcommand{\MatterSpace}{\Psi}

\newcommand{\Car}{\mathrm{car}}
\newcommand{\Cur}{\mathrm{cur}}
\newcommand{\Hom}{\mathrm{hom}}

\newcommand{\ForSpace}[1]{\mathcal I_{#1}^{\mathrm{for}}}
\newcommand{\PrimShell}{\mathcal S_{\mathrm{prim}}}
\newcommand{\MetricOut}{\mathscr G_{\mathrm{out}}}
\newcommand{\MatterOut}{\mathscr T_{\mathrm{out}}}

\newcommand{\ProtoSectionPackage}{\mathfrak Q^{\mathrm{sub,proto,secgen}}}
\newcommand{\ProtoSectionBody}[1]{\mathcal B_{#1}^{\mathrm{psec}}}

\newcommand{\PreProtoSectionPackage}{\mathfrak R^{\mathrm{sub,proto,presec}}}
\newcommand{\PreProtoSectionMinSectionCore}{\mathfrak R^{\mathrm{sub,proto,presec,sec}}}
\newcommand{\PreProtoSectionResidualPackage}{\mathfrak R^{\mathrm{sub,proto,presec,res}}}

\newcommand{\PrePreProtoSectionPullbackCore}{\mathfrak S^{\mathrm{sub,proto,pppresec,pb}}}
\newcommand{\PrePreProtoSectionVertical}[1]{V_{#1}^{\mathrm{pppresec}}}
\newcommand{\PrePreProtoSectionResidualBody}[1]{\mathcal D_{#1}^{\mathrm{pppresec}}}
\newcommand{\PrePreProtoSectionBodyResponse}[1]{\mathscr R_{#1}^{\mathrm{pppresec,pb,body}}}

\newcommand{\PullbackOperatorCore}{\mathfrak S^{\mathrm{sub,proto,pppresec,pb,op}}}
\newcommand{\PullbackBodyOperatorCore}{\mathfrak S^{\mathrm{sub,proto,pppresec,pb,bodyop}}}
\newcommand{\PullbackBodyResponseCore}{\mathfrak S^{\mathrm{sub,proto,pppresec,pb,bodyresp}}}
\newcommand{\RecoveredBodyResponseSpace}[1]{\widehat V_{#1}^{\mathrm{pppresec,pb,bodyresp}}}
\newcommand{\BodyResponseReduction}{\mathfrak B_{\mathrm{bodyresp}}}
\newcommand{\BodyResponseDatum}{\mathfrak D_{\mathrm{bodyresp}}}

\newtheorem{definition}{Definition}
\newtheorem{proposition}{Proposition}
\newtheorem{remark}{Remark}
\newtheorem{corollary}{Corollary}
\newtheorem{theorem}{Theorem}

\title{\textbf{The \TheoryName{}}\\[0.4em]
\large Primitive-Reservoir Observer-Reduction\\
\large Section-Centered Hidden-Response\\
\large Observer-Relevant Pullback Body-Response Core\\
\large Necessity / Uniqueness Theorem}
\author{}
\date{}

\begin{document}

\maketitle

\begin{abstract}
The current active-primary theorem already removed the full observer-relevant pullback body-operator core \(\PullbackBodyOperatorCore\) as the primitive stopping point of the active section-centered hidden-response line. Benchmark-facingly, the separate ambient operator-family presentation is now recovered from the smaller observer-relevant pullback body-response core \(\PullbackBodyResponseCore\). The current benchmark criterion therefore made the honest next move explicit: either derive or justify that body-response core from still more primitive explicit substrate structure, or, if no such clean deeper law is yet in hand, prove a sharper theorem strictly inside that smaller core.

The present note takes that honest internal route. It proves that on the current section-centered hidden-response route the observer-relevant pullback body-response core is already necessary and unique at benchmark scope up to body-response-faithful equivalence. More precisely, any benchmark-faithful reduction that preserves the fixed proto-section benchmark base, the current body-response recovery route, and the benchmark one-metric observer package must determine the same section-centered displacement-body datum together with the same body-restricted pullback response. Conversely, nothing beyond that core is needed at current benchmark scope, because the earlier pullback-body-response-core collapse theorem already shows that all present downstream benchmark-facing uses factor through that smaller core.

This is not yet a completed first-principles derivation of GR from substrate microphysics. Its positive result is narrower and more honest. It removes the remaining presentational freedom at the current body-response frontier on the recorded route. The next primary manuscript must therefore derive or uniquely force the current body-response core itself, or some nontrivial constituent or relation inside it, from still more primitive explicit substrate structure. Another same-output deeper relay or a re-expansion back to the larger pullback-body-operator, pullback-operator, pullback-quadratic, hidden-response, residual-normal-form, or pre-proto-section presentations is no longer the right front-line move.
\end{abstract}

\section{Introduction}

The active derivational program of \TheoryName{} now has a sharper task than merely going deeper again. The current pullback-body-response-core collapse theorem already spent the larger observer-relevant pullback body-operator core as the primitive stopping point on the active section-centered hidden-response route. Once the actual displacement body together with the quadratic pullback response realized on that body are fixed, the full ambient operator-family presentation is recovered and all current downstream benchmark-facing consequences follow unchanged. The present frontier is therefore already internal to the body-response core itself.

That status leaves only two honest options. One can derive or justify the current body-response core from still more primitive explicit substrate structure in a way that genuinely changes benchmark-facing status. Or, if that deeper forcing law is not yet cleanly available, one can prove a sharper internal theorem that removes remaining presentational freedom strictly inside the current body-response core. The present note takes the second route. Its claim is that, on the currently recorded route, the pullback body-response core is already the unique minimal benchmark-facing datum up to body-response-faithful equivalence.

\paragraph{Framework and claim boundary.}
\TheoryProgram{} denotes the broader \Aether{} / \Aflow{} framework built on the claims that the \Aether{} is the underlying four-dimensional substrate of reality, the \Aflow{} is its intrinsic ordered motion, observed three-dimensional space is the local experiential slice of that deeper substrate, S-time is the experienced order of change arising from matter, light, and the \Aflow{}, and observed expansion is the three-dimensional appearance of deeper four-dimensional ordered motion. Within that framework, \TheoryName{} denotes the exact-closure theory in which the effective gravitational dynamics are adopted to be exactly Einsteinian with universal matter coupling. In the active sequence, \emph{adoption} means use of that established relativistic dynamics without claiming substrate derivation, while \emph{derivation} is reserved for a first-principles recovery from explicit substrate variables.

The present note stays strictly inside that boundary. It does not introduce a second operative metric, it does not enlarge the low-energy matter law, and it does not reopen the larger pullback-body-operator, pullback-operator, pullback-quadratic, hidden-response, residual-normal-form, or pre-proto-section presentations as if they were again independently live. Its narrower benchmark-facing role is to prove that the current body-response core is already the unique minimal datum of the recorded route at current scope.

\section{Fixed Inputs from the Current Body-Response Frontier}

\begin{definition}[Recorded primitive continuation data]
\label{def:fixed}
Fix the following continuation data from the active primitive-reservoir observer-reduction line.
\begin{enumerate}
    \item For each sector label \(\sigma\in\{\Car,\Cur,\Hom\}\), a primitive forcing shell
    \[
    \ForSpace{\sigma}.
    \]
    \item The product shell
    \[
    \PrimShell
    \equiv
    \ForSpace{\Car}\times \ForSpace{\Cur}\times \ForSpace{\Hom}.
    \]
    \item The recorded metric and ordinary-matter outputs
    \[
    \MetricOut:\PrimShell\to\{\text{observer metrics}\},
    \qquad
    \MatterOut:\PrimShell\times\MatterSpace\to\{\text{observer stress tensors}\},
    \]
    satisfying
    \begin{equation}
    \MetricOut(\mathbf w)=\CompletedMetric,
    \qquad
    \MatterOut(\mathbf w,\psi)
    =
    T_{\mu\nu}^{\mathrm{matter}}[\CompletedMetric,\psi]
    \label{eq:fixedoutputs}
    \end{equation}
    for every \(\mathbf w\in\PrimShell\) and every matter configuration \(\psi\in\MatterSpace\).
\end{enumerate}
\end{definition}

\begin{definition}[Recorded proto-section benchmark base and current route]
\label{def:route}
Fix \(\sigma\in\{\Car,\Cur,\Hom\}\). The active route already fixes the following data and consequences.
\begin{enumerate}
    \item the current proto-section benchmark base \(\ProtoSectionPackage\), in particular the recorded section body \(\ProtoSectionBody{\sigma}\);
    \item the current observer-relevant pullback body-response core \(\PullbackBodyResponseCore\), the current observer-relevant pullback body-operator core \(\PullbackBodyOperatorCore\), the current observer-relevant pullback-operator core \(\PullbackOperatorCore\), the current observer-relevant pullback-quadratic core \(\PrePreProtoSectionPullbackCore\), the current section-centered residual-normal-form realization \(\PreProtoSectionResidualPackage\), the current observer-relevant pre-proto-section minimizer-section core \(\PreProtoSectionMinSectionCore\), and the current benchmark-faithful pre-proto-section package \(\PreProtoSectionPackage\);
    \item the previously recorded fact that, once a benchmark-faithful observer-relevant pullback body-response core over the fixed proto-section benchmark base is available, the current body-operator-core realization, the current pullback-operator-core realization, the current pullback-quadratic-core realization, the current section-centered residual-normal-form realization, the current pre-proto-section package, the current minimizer-section core, and the previously recorded seed, root, foundation, ground, origin, precursor, lift, shell-restriction, split-core, selector-law, combined-response, ambient-package, trace-core, exact-shell-affine-nucleus, continuity-Hessian compressed core, and benchmark one-metric observer-package consequences follow unchanged;
    \item benchmark faithfulness, meaning preservation of the benchmark outputs \eqref{eq:fixedoutputs}, no second operative metric, no enlarged low-energy matter law, no change to the ordinary matter route, and no premature promotion from adoption to completed derivation.
\end{enumerate}
\end{definition}

\begin{definition}[Observer-relevant pullback body-response core]
\label{def:bodyresp}
Fix \(\sigma\in\{\Car,\Cur,\Hom\}\). A \emph{benchmark-faithful observer-relevant pullback body-response core} \(\PullbackBodyResponseCore\) for sector \(\sigma\) consists of:
\begin{enumerate}
    \item a compact convex section-centered displacement body
    \[
    \PrePreProtoSectionResidualBody{\sigma}
    \subset
    \ProtoSectionBody{\sigma}\times\PrePreProtoSectionVertical{\sigma}
    \]
    whose fiber
    \[
    \PrePreProtoSectionResidualBody{\sigma}(y)
    \equiv
    \left\{
    v\in\PrePreProtoSectionVertical{\sigma}:(y,v)\in\PrePreProtoSectionResidualBody{\sigma}
    \right\}
    \]
    is nonempty, compact, convex, contains \(0\), and has nonempty interior for every \(y\in\ProtoSectionBody{\sigma}\);
    \item the body-restricted pullback response
    \[
    \PrePreProtoSectionBodyResponse{\sigma}:
    \PrePreProtoSectionResidualBody{\sigma}\to[0,\infty),
    \]
    smooth in \(y\), fiberwise quadratic in \(v\), and satisfying
    \[
    \PrePreProtoSectionBodyResponse{\sigma}(y,v)=0
    \qquad\Longleftrightarrow\qquad
    v=0
    \]
    on every fiber;
    \item benchmark faithfulness in the sense of Definition \ref{def:route}(4).
\end{enumerate}
From this core define the recovered ambient displacement space by
\[
\RecoveredBodyResponseSpace{\sigma}
\equiv
\operatorname{span}\left(
\bigcup_{y\in\ProtoSectionBody{\sigma}}
\PrePreProtoSectionResidualBody{\sigma}(y)
\right)
\subset
\PrePreProtoSectionVertical{\sigma}.
\]
\end{definition}

\begin{remark}
Definition \ref{def:bodyresp} is the current minimal data package identified by the preceding collapse theorem. The present note asks whether, on the recorded route itself, any smaller genuinely benchmark-facing datum still remains below it.
\end{remark}

\begin{proposition}[What the preceding collapse theorem already proves]
\label{prop:collapse}
At current scope, the earlier pullback-body-response-core collapse theorem already proves all of the following.
\begin{enumerate}
    \item The ambient displacement space is recovered canonically from the fiberwise span of the displacement body.
    \item The body-restricted pullback response determines the unique compatible positive self-adjoint pullback operator family on that recovered space.
    \item The full pullback body-operator core is reconstructed from the pullback body-response core together with the fixed proto-section benchmark base.
    \item All current benchmark-facing downstream uses of the body-response frontier factor through \(\PullbackBodyResponseCore\).
\end{enumerate}
\end{proposition}

\begin{proof}
This is exactly the content of the earlier theorem collapsing the observer-relevant pullback body-operator core to the observer-relevant pullback body-response core and of its downstream factorization statement.
\end{proof}

\section{Body-Response-Faithful Benchmark Reductions}

\begin{definition}[Body-response-faithful benchmark reduction]
\label{def:reduction}
A \emph{body-response-faithful benchmark reduction} of the current section-centered hidden-response route is a reduction \(\BodyResponseReduction\) from deeper explicit substrate variables to the observer-facing sector such that, after admissible elimination of auxiliary or substrate-only variables and for each \(\sigma\in\{\Car,\Cur,\Hom\}\),
\begin{enumerate}
    \item it preserves the fixed proto-section benchmark base of Definition \ref{def:route}(1);
    \item it determines a benchmark-faithful candidate body-response datum
    \[
    \BodyResponseDatum(\BodyResponseReduction,\sigma)
    =
    \bigl(
    \widetilde{\mathcal D}_{\sigma},
    \widetilde{\mathscr R}_{\sigma}
    \bigr)
    \]
    of the same form as Definition \ref{def:bodyresp};
    \item it preserves the current route in the precise benchmark-facing sense that the candidate datum reconstructs a benchmark-faithful pullback body-operator core and hence the same current pullback-operator core, pullback-quadratic core, section-centered residual-normal-form realization, pre-proto-section package, continuity-Hessian compressed core, and benchmark one-metric observer package;
    \item it introduces no second operative metric, no enlarged low-energy matter law, and no change to the ordinary matter route.
\end{enumerate}
\end{definition}

\begin{remark}
Definition \ref{def:reduction} is intentionally route specific. The theorem of this note concerns the active recorded section-centered hidden-response route, not every imaginable future redesign that changes the route itself.
\end{remark}

\begin{definition}[Body-response-faithful equivalence]
\label{def:equiv}
Fix \(\sigma\in\{\Car,\Cur,\Hom\}\).
Two candidate body-response data
\[
\bigl(\mathcal D_{\sigma}^{(1)},\mathscr R_{\sigma}^{(1)}\bigr),
\qquad
\bigl(\mathcal D_{\sigma}^{(2)},\mathscr R_{\sigma}^{(2)}\bigr)
\]
over the fixed proto-section benchmark base are \emph{body-response-faithfully equivalent} if, for every \(y\in\ProtoSectionBody{\sigma}\), there is a linear isomorphism
\[
\Phi_{\sigma}(y):
\operatorname{span}\bigl(\mathcal D_{\sigma}^{(1)}(y)\bigr)
\xrightarrow{\cong}
\operatorname{span}\bigl(\mathcal D_{\sigma}^{(2)}(y)\bigr)
\]
such that
\[
\Phi_{\sigma}(y)\bigl(\mathcal D_{\sigma}^{(1)}(y)\bigr)
=
\mathcal D_{\sigma}^{(2)}(y)
\]
and
\[
\mathscr R_{\sigma}^{(2)}\bigl(y,\Phi_{\sigma}(y)v\bigr)
=
\mathscr R_{\sigma}^{(1)}(y,v)
\qquad
\text{for every }
v\in\mathcal D_{\sigma}^{(1)}(y).
\]
\end{definition}

\begin{definition}[Minimal observer-relevant body-response datum]
\label{def:minimal}
For a body-response-faithful benchmark reduction \(\BodyResponseReduction\), a datum \(\BodyResponseDatum(\BodyResponseReduction)\) is a \emph{minimal observer-relevant body-response datum} if:
\begin{enumerate}
    \item it suffices to determine every benchmark-facing derivation statement available on the current section-centered hidden-response route;
    \item no proper subdatum of it still suffices at that same benchmark-facing scope.
\end{enumerate}
\end{definition}

\section{Necessity of the Current Body-Response Constituents}

\begin{proposition}[The displacement-body datum is necessary]
\label{prop:body}
Every minimal observer-relevant body-response datum for a body-response-faithful benchmark reduction must determine the section-centered displacement-body datum on the current route up to body-response-faithful equivalence.
\end{proposition}

\begin{proof}
The current route does not use an abstract quadratic response with unspecified domain. It uses the recorded section-centered displacement body on which the response is actually realized. By Definition \ref{def:bodyresp}, that body records the admissible displacement fibers, contains the zero section, and has nonempty interior in each fiber. Proposition \ref{prop:collapse}(1) then uses that body to recover the ambient displacement space itself. If the displacement-body datum were omitted, neither the domain of the current body response nor the recovered ambient-space step of the current route would be fixed. Therefore no minimal benchmark-facing datum on this route may omit the displacement-body constituent.
\end{proof}

\begin{proposition}[The body-restricted pullback response is necessary]
\label{prop:response}
Every minimal observer-relevant body-response datum for a body-response-faithful benchmark reduction must determine the body-restricted pullback response on the current route up to body-response-faithful equivalence.
\end{proposition}

\begin{proof}
The current route does not recover the pullback body-operator core, the pullback-operator core, or the remaining downstream benchmark-facing package from the displacement body alone. Proposition \ref{prop:collapse}(2) uses the body-restricted quadratic response to recover the unique compatible pullback operator family, and Proposition \ref{prop:collapse}(4) then carries all current downstream benchmark-facing uses through that recovery. If the response constituent were omitted, the present route would no longer determine the operator recovery or the downstream package. Therefore no minimal benchmark-facing datum on this route may omit the response constituent.
\end{proof}

\begin{proposition}[The current body-response core is sufficient]
\label{prop:sufficient}
The current observer-relevant pullback body-response core \(\PullbackBodyResponseCore\) is sufficient to determine every benchmark-facing derivation statement available on the current section-centered hidden-response route.
\end{proposition}

\begin{proof}
By Proposition \ref{prop:collapse}(4), all current benchmark-facing downstream uses of the present frontier already factor through \(\PullbackBodyResponseCore\). Therefore the current body-response core is sufficient at benchmark scope on the recorded route.
\end{proof}

\section{Main Theorem}

\begin{theorem}[Observer-relevant pullback body-response core necessity / uniqueness on the current route]
\label{thm:main}
For every body-response-faithful benchmark reduction \(\BodyResponseReduction\) of the current section-centered hidden-response route, the observer-relevant pullback body-response core \(\PullbackBodyResponseCore\) of Definition \ref{def:bodyresp} is necessary and unique at benchmark scope up to body-response-faithful equivalence. More precisely:
\begin{enumerate}
    \item every minimal observer-relevant body-response datum \(\BodyResponseDatum(\BodyResponseReduction)\) must determine both the displacement-body datum and the body-restricted pullback response on the current route;
    \item \(\PullbackBodyResponseCore\) itself is sufficient to determine every benchmark-facing derivation statement available on that route;
    \item any two minimal observer-relevant body-response data are body-response-faithfully equivalent to \(\PullbackBodyResponseCore\), and hence to each other.
\end{enumerate}
\end{theorem}

\begin{proof}
Item (1) is the combined content of Propositions \ref{prop:body} and \ref{prop:response}. Those propositions show that no minimal observer-relevant datum on the current route may omit either the displacement-body constituent or the body-restricted pullback-response constituent.

Item (2) is Proposition \ref{prop:sufficient}. The earlier collapse theorem already proves that all current benchmark-facing downstream uses factor through \(\PullbackBodyResponseCore\).

For item (3), let \(\BodyResponseDatum^{(1)}\) and \(\BodyResponseDatum^{(2)}\) be two minimal observer-relevant body-response data for body-response-faithful benchmark reductions of the current route. By item (1), each must determine the same kind of displacement-body and body-response constituents through which the route recovers the benchmark-facing downstream package. By item (2), neither needs anything beyond those constituents at current benchmark scope. Therefore each is body-response-faithfully equivalent to \(\PullbackBodyResponseCore\) in the sense of Definition \ref{def:equiv}. Hence they are body-response-faithfully equivalent to each other as well.
\end{proof}

\begin{corollary}[No smaller benchmark-facing datum remains on the current route]
\label{cor:minimal}
On the current section-centered hidden-response route, there is no genuinely smaller benchmark-facing datum below \(\PullbackBodyResponseCore\) up to body-response-faithful equivalence.
\end{corollary}

\begin{proof}
If a smaller datum existed, it would either omit the displacement-body constituent or the body-restricted pullback-response constituent proved necessary by Propositions \ref{prop:body} and \ref{prop:response}, or else merely rename the same constituents and therefore remain body-response-faithfully equivalent to \(\PullbackBodyResponseCore\). The first option contradicts necessity, and the second is not genuinely smaller at benchmark scope.
\end{proof}

\begin{remark}
Corollary \ref{cor:minimal} is the benchmark-facing sense in which the current body-response frontier is now exhausted internally on the recorded route. The route no longer has live benchmark-facing freedom in another body-response-core restatement or in a re-expansion to a larger presentation.
\end{remark}

\section{What Must Be Derived Next}

\begin{corollary}[Primary post-uniqueness burden]
\label{cor:next}
After Theorem \ref{thm:main}, the next primary manuscript below the current frontier must do at least one of the following:
\begin{enumerate}
    \item derive the current displacement-body datum from still more primitive explicit substrate structure rather than preserving it as already-recorded route data;
    \item derive the current body-restricted pullback response from still more primitive explicit substrate structure rather than preserving it as already-recorded route data;
    \item prove a sharper necessity, uniqueness, or compatibility theorem for some nontrivial constituent or relation inside the present body-response core;
    \item do one of the previous three while preserving the same benchmark one-metric package, the same ordinary matter route, and the same claim boundary.
\end{enumerate}
Another same-output deeper relay that merely preserves the same body-response core, or any re-expansion back to the larger pullback-body-operator, pullback-operator, pullback-quadratic, hidden-response, residual-normal-form, or pre-proto-section presentations, is not primary by default.
\end{corollary}

\begin{proof}
By Theorem \ref{thm:main} and Corollary \ref{cor:minimal}, the current body-response core is already the unique minimal benchmark-facing datum on the recorded route up to body-response-faithful equivalence. Therefore a new primary manuscript must derive or sharpen some constituent or relation inside that datum. If it does not, then it preserves the same benchmark-facing core and is neutral at project-status level.
\end{proof}

\begin{proposition}[The benchmark package remains fixed]
\label{prop:benchmark}
The present necessity / uniqueness theorem preserves the benchmark one-metric package \eqref{eq:fixedoutputs}. In particular, the operative metric remains exactly \(\CompletedMetric\), the ordinary matter route is unchanged, there is no second operative metric, and the low-energy matter law is not enlarged.
\end{proposition}

\begin{proof}
This is part of benchmark faithfulness in Definitions \ref{def:route}(4), \ref{def:bodyresp}(3), and \ref{def:reduction}(4). The present note removes only the remaining presentational freedom at current body-response scope on the recorded route.
\end{proof}

\section{Conclusion}

The positive result of this note is narrower than a completed substrate derivation of GR but sharper than another routing reminder. On the current section-centered hidden-response route, the observer-relevant pullback body-response core is now proved to be necessary and unique at benchmark scope up to body-response-faithful equivalence.

That core consists of the section-centered displacement-body datum together with the body-restricted pullback response realized on that body, under the same benchmark-faithfulness boundary. By the theorem of this note, no smaller genuinely benchmark-facing datum remains below it on the recorded route. The current pullback-body-response-core collapse theorem remains preserved main-line history because it first moved the live burden below the larger body-operator-core presentation. The present necessity / uniqueness theorem then spends the remaining presentational freedom at that smaller frontier.

This sharpens the honest next burden. A new primary manuscript must no longer merely show that an even deeper substrate layer reproduces the same body-response core or re-expands to a larger already-spent presentation. It must derive or uniquely force the current body-response core itself, or some nontrivial constituent or relation inside it, from still more primitive explicit substrate structure while preserving the same benchmark one-metric package, the same ordinary matter route, and the same claim boundary.

\end{document}
